\chapter{The Precedents and the Playbook}
\label{c:precedents}

\begin{glance}
Solar showed the pure form \full{3--8}: Western invention and demand, Chinese policy-financed scale, capture of the learning curve, failed tariffs, over 90\% share at every upstream stage, then state cartelization. Batteries added vertical depth from mine to pack, won by regulatory market reservation and by industrializing the ``inferior'' chemistry until the West licensed it back. Electric vehicles proved the playbook works on a whole system. Telecommunications, a fourth and lesser precedent, shows what it does with a financed network---and supplies the template the sovereign-AI-export story of Chapter~\ref{c:world} is borrowed from. From the campaigns the book distills nine tactical moves and six strategic steps; from those that have \emph{not} worked, where it grinds: tacit knowledge, slow clock speed, foreign certification, installed-base lock, trust. AI's model and compute layers sit in the kill zone; its deployment layer is genuinely contested.
\end{glance}

\section{Three campaigns, in brief}

\textbf{Solar} was invented at Bell Labs and built on American, German, Japanese and Australian science; Sharp was the largest cell producer in 2006, and demand was created by Germany's feed-in tariff. Chinese scale-up rode that demand and about \$47 billion of development-bank credit~\cite{cdb-credit}; polysilicon fell from about \$475 to under \$16 per kilogram between 2008 and 2012~\cite{bernreuter}, and Solyndra, Q-Cells and SolarWorld died---as did the Chinese pioneer Suntech, because the state protects the industry, not the champion. US and EU duties from 2012 taxed deployment without resurrecting manufacturing, and a domestic market of 315\,GW$_{AC}$ in 2025 alone made foreign trade policy irrelevant~\cite{china-2025-installs}. The end state is over 90\% Chinese share at every stage upstream of the module~\cite{pv-shares-2023}, module prices down from \$106 per watt in 1976 to about \$0.09 at a learning rate near 20\% per doubling~\cite{owid-learning,energytrend-prices}---and since roughly 2010 essentially all doublings were Chinese, so whoever manufactures the doublings owns the learning. Winning is expensive: losses below cash cost, then Xi's July 2025 anti-involution campaign and a polysilicon buyout platform, an OPEC for solar~\cite{anti-involution,poly-platform}. Four lessons: trade barriers after cost-curve capture protect only domestic prices; champions are disposable, the industrial commons is not; a home market absorbing half of world demand makes foreign trade policy irrelevant; and commoditization ends in state-managed consolidation. Keep them in mind when reading about open-weight models priced near zero.

\textbf{Batteries} were commercialized by Sony in 1991 and led by Japan and Korea for two decades; China entered through low-margin consumer cells and pivoted when the state created demand. The decisive instrument was a certification, not a tariff: the 2015 white list of 57 domestic makers eligible for NEV subsidies, which excluded LG and Samsung and was scrapped in 2019 when no longer needed~\cite{white-list}. By 2025 Chinese makers held about 69\% of installations and over 80\% of the world's 4\,TWh of cell capacity~\cite{sne-2025,iea-gevo-2026}, plus 70--90\% of most upstream layers and two thirds of lithium and cobalt refining, at pack costs that leave North America and Europe paying a 44--56\% premium~\cite{bnef-2025}: a gigafactory built outside China is built, in effect, inside the Chinese supply chain. Lithium iron phosphate---Western-invented, dismissed as a dead end---passed 55\% of the global EV market and is licensed back to Ford and Stellantis, the inferior branch industrialized at scale beating the premium branch on system economics, a preview of the LPDDR-versus-HBM argument of Chapter~\ref{c:elements}. In 2025 the infant industry became the gatekeeper with export controls on battery technology~\cite{mofcom-battery-controls}. Three more lessons: regulatory reservation beats tariffs and is invisible to trade law; vertical integration to the mine makes a foreign factory a hostage; watch what China industrializes, not what the West benchmarks.

\textbf{Electric vehicles} are a whole system wrapped in a century of incumbency; if the playbook worked here, no consumer-industrial market is safe, and it worked. Wan Gang's leapfrog decision predates Tesla; at least \$231 billion of layered demand support~\cite{csis-230b}, a dual-credit mandate binding foreign joint ventures, and 19 million charging points took NEV penetration from 6\% in 2020 to 54\% in 2025~\cite{cpca-2025} and made China the largest car exporter~\cite{caam-2025}. BYD ran the flywheel from 416,000 to 4.55 million vehicles with about 80\% of Tier-1 content in-house and a \$4,700 structural advantage over Tesla of which subsidy is under \$300~\cite{rhodium-evs}, while some 400 Chinese EV ventures died around it: the state subsidizes the species, not the specimen. Foreign brands fell from 62\% to 35\% of the Chinese market and technology now flows backward. And the car became a software platform in Chinese hands---about 95\% of automotive lidar, driver assistance free on an \$8,000 car, 20--24-month cycles against 40--50---so the first mass-market safety-critical embedded-AI category is Chinese-led, and it anchors the domestic chips, sensors and models of the next chapter. Three final lessons: pick the transition where incumbency resets to zero and force it at home; a system-level win pulls the whole component stack with it; speed compounds.

\textbf{Telecommunications, in lesser detail}, is the precedent closest in structure to what an ``AI stack'' export would be---infrastructure, financed, standardized, locked in by service---and one company ran almost all of the playbook on it. Huawei held a fifth of the Chinese switch market in 1996 and reached parity with Ericsson, the incumbent it had been imitating, by 2011~\cite{ahrens-huawei}; its share of world telecom-equipment revenue ran 27.7\% in 2018, 31\% in the first half of 2020, the year the sanctions bit hardest, and \textbf{31\% again in the first half of 2025}, 40\% outside North America~\cite{delloro-huawei}. The instrument that distinguishes the campaign, and that AI has not reproduced, was credit: a China Development Bank facility of \$10 billion for Huawei's customers in 2004, raised to \textbf{\$30 billion in December 2009}, inside up to \$75 billion of state support by the \emph{Wall Street Journal}'s reconstruction, buying prices about 30\% below rivals' where the credit was and not elsewhere~\cite{wsj-huawei-state}---an operator in Lagos was choosing between a switch it could finance and one it could not. Entry ran rural-first, then Africa, Asia and Latin America, then Europe; the endgame was standards from inside, first in declared 5G patent families and in 3GPP contributions~\cite{huawei-5g-seps}. The bans---US Entity List 2019, Britain's reversal 2020, an EU phase-out proposal in 2026---removed share only where security agencies held a veto and operators could swap at no cost, the Five Eyes, the Nordics and Japan: Chinese vendors held about 32\% of European 5G sites in 2022 and about 32\% at the end of 2024, and what actually turned Britain was the May 2020 \emph{foundry} rule, a semiconductor chokepoint, not a network-security finding~\cite{huawei-bans}. The disanalogy, tested in Chapter~\ref{c:world}: AI has no credit line, no surplus capacity, and an artifact---open weights---that diffuses better than any switch and locks in nothing \full{6}.

\section{The playbook distilled}

Table~\ref{tab:playbook} inventories nine tactical moves that recur across the campaigns; the AI column is developed in Chapter~\ref{c:elements}. Six strategic steps are the sequence the moves serve. \emph{Target selection}: a discontinuity where the incumbent's accumulated advantage is about to be stranded and where manufacturing learning-by-doing, not protected IP, is the accumulating asset. \emph{Demand before supply}: the home market as guaranteed demand, foreign firms admitted just long enough to transfer capability. \emph{Overcapacity as strategy}: capacity at multiples of world demand, which Western economics reads as waste and which purchases the entire learning curve; the losses land on Chinese balance sheets and dead firms while the cost curve remains. \emph{Selection, not planning}: Beijing floods the sector and lets domestic combat select, its loyalty being to the industrial commons---the engineer pool, the supplier mesh, the tooling ecology. \emph{Commoditize the adversary's rent layer}: wherever the Western business model collects rent---module, pack, badge, API token, ISA licence---price it at marginal cost, because rents fund the West's R\&D flywheel, and openness is this move in its purest form. \emph{Endgame}: consolidate, stabilize prices upward, and convert dominance into coercive leverage.

\begin{table}[htbp]
\centering\footnotesize
\caption{The nine tactical moves, across four campaigns; telecommunications is scored against the six steps in the full edition rather than tabulated here.}
\label{tab:playbook}
\begin{tabular}{@{}p{26mm}p{25mm}p{25mm}p{25mm}p{40mm}@{}}
\toprule
\textbf{Move} & \textbf{Solar} & \textbf{Batteries} & \textbf{EVs} & \textbf{AI} \\
\midrule
M1 State designation & SEI 2010, 12th FYP & White list 2015 & 863 program, SEI 2010 & ``AI+'' plan, 15th FYP, RISC-V guidance \\
M2 Home demand & National FiT, Top Runner & NEV subsidy linkage & \$231B support, dual credit & State DCs mandated on domestic chips, subsidized power \\
M3 Subsidized scale-up & CDB \$47B credit & Provincial capital & Provincial EV funds & Big Fund III (\$47.5B), SMIC/CXMT capex \\
M4 Darwinian competition & 2018 ``531'' cull & 100+ makers to consolidation & $\sim$400 dead EV firms & LLM price war since 2024; chip-vendor shakeout beginning \\
M5 Supply-chain capture & poly $\to$ module & mine $\to$ cell & battery $\to$ chip $\to$ magnet $\to$ ship & model $\to$ framework $\to$ chip $\to$ interconnect $\to$ fab $\to$ DRAM \\
M6 Commoditize the rent layer & module \$/W & pack \$/kWh & vehicle BOM & open weights, \$/Mtoken, open ISA and stack \\
M7 Absorb the ``inferior'' branch & mono-PERC, TOPCon & LFP & small BEVs & sparse MoE; LPDDR capacity vs HBM bandwidth \\
M8 Standards capture & efficiency floors & GB safety standards & charging standards & RISC-V profiles, matrix ISA, UnifiedBus \\
M9 Leverage after victory & OPEC-for-poly & tech export controls 2025 & rare-earth controls 2025 & (pending) \\
\bottomrule
\end{tabular}
\end{table}

Against this the West has fielded two counter-patterns, both failing. Tariff walls arrive after cost-curve capture and merely tax domestic consumers while the walled market shrinks in relevance. Chokepoint denial---semiconductor export controls since 2019---is the more serious strategy, and its record is a pattern the book calls \emph{chokepoint decay}: each denied input (EUV, HBM, H100s, EDA) triggers a state-scale substitution program plus an architectural detour around the input, while the denial hands the domestic substitute a guaranteed market. Export controls are, functionally, a US-funded market-reservation policy for Huawei, SMIC, CXMT and Cambricon.

\section{Where the analogy breaks}
\label{c:analogy}

A playbook argued from three victories owes the reader its defeats. The same state, spending an estimated \$248 billion a year on industrial policy~\cite{red-ink}, has run two-decade campaigns where dominance has not arrived: commercial aircraft (COMAC's C919 delivered about thirteen aircraft in 2025 against a plan of 75, on a Western engine whose export licences Washington suspended and restored---chokepoint leverage working); lithography and the deep tool chain (under 5\% localized after twenty years and three Big Funds~\cite{litho-localization,litho}); precision machine tools (five-axis CNC about 8\% domestic against an 80\% target); operating systems, where HarmonyOS reached 17\% only after sanctions burned the Android bridge---ecosystems yield on decade timescales, under existential forcing, inside the home market; and biotech, where \$138 billion of 2025 out-licensing meets a residual gap that is entirely trust and regulation, the FDA's 14--1 rejection of China-only trial data being the certification barrier in its purest form~\cite{biotech}. The playbook wins where products are modular, iteration is fast, learning is manufacturable and the home market can certify; it grinds where knowledge is tacit, certification is foreign, and value lives in an installed ecosystem.

Scored against that model, AI models are the most modular, fastest-clock-speed artifact in industrial history---sparse-MoE training is closer to batteries than to jet engines---and compute is intermediate, with the tool chain inheriting lithography's resistance and being routed around. But three properties of AI sit in the resisting column and are the strongest structural arguments against this book's title: the frontier moves, and recursive AI-for-AI research could reopen the gap faster than any factory closes it; ecosystems lock, and the enterprise deployment layer is Office, not modules; and trust and certification bind, with the sintilimab precedent---capability accepted, certification denied---directly available to Western regulators. Two asymmetries follow: zero reproduction cost lets Chinese weights conquer adoption instantly and would let a superior Western open release reconquer it just as fast, so there is no factory-shaped moat around P1; and talent is the input the two systems still share---researchers of Chinese undergraduate origin were 47\% of the top tier in 2022, up from 29\% in 2019, while the share of that tier working in the United States fell from 59\% to 42\% and China's rose from 11\% to 28\%~\cite{talent}. Model and compute therefore sit mostly in the kill zone, the tool chain in the resistance zone and bypassed, and the deployment layer---trust, certification, ecosystems---is contested ground where the precedents predict a long grind, not a collapse.
